\documentclass[openright, oneside, a4paper, article, brazil]{abntex2}
\usepackage[no-math]{fontspec}
\setmainfont{NewJune}[
	UprightFont = *-Regular,
	BoldFont=*-Bold,
	ItalicFont=*-Regular-Italic,
	BoldItalicFont=*-Bold-Italic,
]
\setsansfont{NewJune}[
	UprightFont = *-Regular,
	BoldFont=*-Bold,
	ItalicFont=*-Regular-Italic,
	BoldItalicFont=*-Bold-Italic,]
\fontsize{12pt}{18pt}
\usepackage{xcolor}
\usepackage%[
	%bookmarks=true, 
	%bookmarksopen=true
	%]
	{hyperref}
\hypersetup{%
	colorlinks=false, 
	allbordercolors=white,
	pdftitle = {Requerimento de cancelamento de inscrição de pessoa fisica},
	pdfauthor = {Conselho Regional de Psicologia de São Paulo},
	pdfsubject = {Cancelamento de registro no CRP SP — pessoa física},
	pdfkeywords = {%
		CRP, 
		inscrição, 
		registro,
		cancelamento, 
		pessoa física}
	}
\usepackage{sectionbreak}
\usepackage[%
    a4paper,
    left = 1.5cm, right = 2cm, top = 2.5cm, bottom = 2cm,
    ]%
{geometry}
\usepackage{graphicx}
\usepackage{microtype}
\usepackage{ragged2e}
\usepackage{xhfill}
\makepagestyle{crp}
\makeoddhead{crp}{}{\includegraphics[height = 2.2cm]{logo-crp.png}}{}
\DoubleSpacing
%\nonzeroparskip
\usepackage{enumitem,amssymb}
\newlist{todolist}{itemize}{2}
\setlist[todolist]{label=$\square$}
\title{Requerimento de cancelamento de inscrição de pessoa física}
\author{Conselho Regional de Psicologia de São Paulo}
%
\begin{document}
\pagestyle{crp}
\vspace*{\baselineskip}
\centering
\Form
\section*{REQUERIMENTO DE CANCELAMENTO DE INSCRIÇÃO DE PESSOA FÍSICA}%
\addcontentsline{toc}{section}{Requerimento de cancelamento de inscrição de pessoa física}
\vspace*{0.618\baselineskip}

\textbf{Exma./Exmo. Presidenta/e do Conselho Regional de Psicologia do Estado de São Paulo (CRP~SP)}
\vspace*{\baselineskip}

\justifying
\noindent
\begin{minipage}[t]{0.995\textwidth}
	Eu, \xrfill[0pt]{0.4pt}\\%
	CRP~SP~nº \rule{8em}{0.4pt} , %
	RG \rule{16em}{0.4pt} , 
	CPF \xrfill[0pt]{0.4pt}\\%
	endereço para correspondência \xrfill[0pt]{0.4pt}\\%
	nº \rule{3em}{0.4pt} , %
	complemento \rule{4em}{0.4pt} , %
	bairro \rule{18em}{0.4pt} , %
	CEP \xrfill[0pt]{0.4pt}\\
	cidade \rule{16em}{0.4pt} , %
	UF \rule{3em}{0.4pt} , %
	tel. residencial ( \rule{1.5em}{0.4pt} ) \xrfill[0pt]{0.4pt}\\% %
	tel. comercial ( \rule{1.5em}{0.4pt} ) \rule{16em}{0.4pt} , %
	tel. celular ( \rule{1.5em}{0.4pt} ) \xrfill[0pt]{0.4pt}\\%
	\vspace{\dp0}%
\end{minipage}%
\begin{minipage}[t]{0.005\textwidth}
	\raggedleft
	,\\
	,\\
	,\\
	,\\
	,\\
	,\\
\end{minipage}
%

\vspace{-0.37\baselineskip}
\noindent \textit{e-mail} \rule{20em}{0.4pt} , venho respeitosamente requerer o cancelamento da inscrição por não estar exercendo a profissão, com\-pro\-me\-ten\-do-me a requerê-la novamente quando retornar às atividades profissionais.

Estou ciente, também, de que o cancelamento da inscrição não quita débitos anteriores com este Conselho, e de que a anuidade do ano em curso será cobrada proporcionalmente, tendo como base o mês anterior a esta solicitação.

\subsection*{ATENÇÃO}%
\addcontentsline{toc}{subsection}{Atenção}
Para deferimento da presente solicitação, faz-se necessário devolver a carteira de identidade profissional (CIP); caso encaminhe pelo correio, o CRP~SP orienta que este este documento seja enviado por meio de carta registrada, por medida de segurança. Quando isto não for possível, justificar a não entrega por um dos motivos abaixo, anexando o boletim de ocorrência correspondente:

\subsubsection*{Motivo para não entregar a CIP:}%
\addcontentsline{toc}{subsubsection}{Motivo para não entregar a CIP}

\begin{todolist}[labelwidth=2em]
\item extravio;
\item roubo.
\end{todolist}

A/o psicóloga/o acima caracterizada/o, que a seguir assina, afirma serem verdadeiros os dados aqui informados.

\vspace*{1.618\baselineskip}
\centering
\rule{21em}{0.4pt}\\
Assinatura

\end{document}